\section{Introduction}
	
	An \emph{eventown} is a family $\mathcal{F}\subset 2^{[n]}$ such that $|A\cap B|$ is even for any $A,B\in\mathcal{F}$. The famous Eventown theorem of Berlekamp \cite{B69},  also proved independently by Graver \cite{G75}, states that if $\mathcal{F}\subset 2^{[n]}$ is an eventown, then $|\mathcal{F}|\leq 2^{\lfloor n/2\rfloor}$. This bound is also the best possible, and a simple construction showing this can be obtained as follows. Say that a family $\mathcal{F}\subset 2^{[n]}$ is \emph{atomic}, if there exist disjoint sets $A_1,\dots,A_d\subset [n]$ such that $\mathcal{F}$ is the family of all sets $F$ satisfying that either $A_i\subset F$ or $A_i\cap F=\emptyset$ for every $i\in [d]$, and $F$ contains no element not covered by the sets $A_i$. 
	The sets $A_1,\dots,A_d$ are called the {\em atoms} of $\mathcal{F}$. 
	Also, let $S(n,\ell)$ be the atomic family for which $d=\lfloor n/\ell\rfloor$ and all $A_i, i\in [d]$ have size exactly $\ell$. Note that $|S(n,\ell)|=2^{\lfloor n/\ell\rfloor}$, and the size of the intersection of any number of sets in $S(n,\ell)$ is divisible by $\ell$. Therefore, the family $S(n,2)$ is an eventown of size $2^{\lfloor n/2\rfloor}$. 
	%
	Moreover, any eventown family can be completed to a maximal one of size $2^{\lfloor n/2\rfloor}$, see e.g. the book of Babai and Frankl \cite{BF92}, which is also a general reference on intersection problems.
	
	In general, one might be tempted to conjecture that the maximal families $\mathcal{F}\subset 2^{[n]}$, whose all pairwise intersections are divisible by $\ell$, have size close to $2^{(1+o(1))n/\ell}$. However, this turns out to be far from the truth. Frankl and Odlyzko \cite{FO83} proved that if there exists a Hadamard matrix of order $4\ell$, then there exists such a family of size $2^{\Omega(n\log \ell/\ell)}$, and this bound is also the best possible up to the constant factor. On the other hand, it follows from a result of Deza, Erd\H{o}s and Frankl \cite{DEF78}, proved also by Frankl and Tokushige \cite{FT16},  that if we consider uniform families, that is, $\mathcal{F}\subset [n]^{(r)}$, then $|\mathcal{F}|\leq \binom{\lfloor n/\ell\rfloor}{ r/\ell}$ if $n$ is sufficiently large given $r$ and $\ell\mid r$. This bound is also the best possible as witnessed by the family $\mathcal{F}=[n]^{(r)}\cap S(n,\ell)$. Let us emphasize that the condition that $n$ must be large compared to $r$ is necessary, otherwise this would contradict the aforementioned construction of Frankl and Odlyzko.
	
	Despite all the above, if we require that the intersection of \emph{any number} of sets in $\mathcal{F}\subset 2^{[n]}$ must have size divisible by $\ell$, then it is not difficult to show that $|\mathcal{F}|\leq 2^{\lfloor n/\ell\rfloor}$ for any $n$ and $\ell$. Moreover, in this case, $\mathcal{F}$ is contained in some isomorphic copy of $S(n,\ell)$ (we say that two families in $2^{[n]}$ are isomorphic if they are equal up to a permutation of $[n]$). In 1983, Frankl and Odlyzko \cite{FO83} asked whether a similar conclusion holds if we only require that the intersection of any $k$ distinct sets in $\mathcal{F}$ has size divisible by $\ell$, where $k$ is some constant only depending on $\ell$. More precisely, they conjectured that for some $k$, we must have $|\mathcal{F}|\leq 2^{(1+o(1))n/\ell}$ for such a family $\mathcal{F}$. Until recently, it was not even known if the bound $2^{O(n\log \ell/\ell)}$ can be improved for any constant $k$. Indeed, while there are many tools to handle pairwise intersections as they correspond to the scalar product of characteristic vectors, $k$-wise intersections are usually harder to analyse, see, e.g.,  \cite{FS,GS02,MST,SV18,SV,V99} for related results. Also, it was shown in \cite{SV18} that if the conjecture is true, $k$ must depend on $\ell$. In particular, if $\ell$ is a power of $2$, there exist families $\mathcal{F}\subset 2^{[n]}$ such that the intersection of any $k$ sets in $\mathcal{F}$ has size divisible by $\ell$, and $|\mathcal{F}|\geq 2^{c_{k}n\log \ell/\ell}$, where $c_{k}>0$ is a constant only depending on $k$. In this paper we resolve the conjecture of Frankl and Odlyzko in the following strong form.
	
	\begin{theorem}\label{thm:main0}
		Let $\ell$ be a positive integer, then there exists $k=k(\ell)$ such that for every positive integer $n$ the following holds. Let $\mathcal{F}\subset 2^{[n]}$ such that the intersection of any $k$ distinct elements of $\mathcal{F}$ is divisible by $\ell$. Then $|\mathcal{F}|\leq 2^{\lfloor n/\ell\rfloor}+c$, where $c=c(\ell,k)$ is a constant, and $c=0$ if $\ell\mid n$ and $n$ is sufficiently large.
	\end{theorem}
	
	Note that the error term $c$ is needed if $\ell\nmid n$. Indeed, in this case $S(n,\ell)$ is not extremal, one can add a constant number of sets contained in the nonempty set not covered by members of $S(n,\ell)$ while retaining the property that the intersection of every $k$ distinct sets has size divisible by $\ell$.
